\documentclass[aip,jcp,preprint,nobibnotes,nofootinbib]{revtex4-1}

\usepackage[utf8]{inputenc}
\usepackage{amsmath}
\usepackage{amssymb}
\usepackage{array}
\usepackage{braket}
\usepackage[capitalize]{cleveref}

\newcommand{\bvec}[1]{\mathbf{#1}}
\newcommand{\Id}{\bvec{1}}
\newcommand{\code}[1]{\texttt{\detokenize{#1}}}

\begin{document}

\title{Technical Notes: Symmetry in Hartree--Fock Calculations}
\author{Forte2 Developers}
\noaffiliation
\date{\today}

\begin{abstract}
This note describes how Forte2 represents molecular symmetry in atomic-orbital
and molecular-orbital spaces, preserves it during Hartree--Fock (HF)
diagonalization, and selects occupations for a requested electronic symmetry
or orbital configuration. Spatial HF uses one-dimensional irreducible
representations of Abelian point groups. Spin--orbit generalized HF uses their
full double groups, with character projectors for spinor spaces and Slater
determinant projectors for the total electronic state. The derivations specify
the matrix conventions, occupation algorithms, representability restrictions,
and interfaces between the symmetry and SCF components.
\end{abstract}

\maketitle

\section{Scope and notation}
\label{sec:scope}

With \code{System(symmetry=True)}, Forte2 centers the molecule at its center of
mass, reorients it along its principal axes of inertia, and detects a supported
Abelian subgroup. The supported spatial groups are $C_1$, $C_i$, $C_2$, $C_s$,
$C_{2h}$, $D_2$, $C_{2v}$, and $D_{2h}$. All spatial operations used below are
defined in this principal frame. Larger molecular groups are represented by
the detected subgroup. With symmetry disabled the spatial group is $C_1$.

Let $\chi_\mu$ denote an atomic orbital (AO), $\bvec{S}$ the AO overlap matrix,
and $\bvec{C}$ the matrix whose columns contain molecular-orbital (MO)
coefficients:
\begin{equation}
 \varphi_p=\sum_\mu\chi_\mu C_{\mu p},\qquad
 S_{\mu\nu}=\braket{\chi_\mu|\chi_\nu},\qquad
 \bvec{C}^\dagger\bvec{S}\bvec{C}=\Id.
 \label{eq:mo}
\end{equation}
The canonical AO orthogonalizer $\bvec{X}$ spans the retained orbital space and
satisfies $\bvec{X}^\dagger\bvec{S}\bvec{X}=\Id$. It may be rectangular when
near-linear dependencies have been removed. An irrep is denoted by $\Gamma$,
its dimension by $d_\Gamma$, and its character by $\chi_\Gamma(g)$. The identity
irrep has index zero.

The public option \code{target_symmetry} specifies the total electronic
determinant's irrep. The option \code{irrep_occupations} specifies the numbers
of occupied orbitals in each irrep, with the conventions in
\cref{tab:occupations}. These are distinct constraints: several orbital
configurations may produce the same total irrep. Both may be supplied if they
are compatible.

\begin{table}[ht]
\setlength{\tabcolsep}{10pt}
\caption{Meaning of each value in \code{irrep_occupations}. Unlisted irreps have
zero occupation; the dictionary specifies a complete configuration.}
\label{tab:occupations}
\begin{tabular}{@{}lll@{}}
\hline
Method & Count per irrep & Required total \\
\hline
RHF & Doubly occupied spatial orbitals & $N/2$ \\
ROHF, UHF, CUHF & $(n_\alpha,n_\beta)$ spatial orbitals & $(N_\alpha,N_\beta)$ \\
GHF & Singly occupied spinors & $N$ \\
\hline
\end{tabular}
\end{table}

Labels are case insensitive, and integer indices are also accepted. Spatial
indices use Cotton ordering. Double-group bosonic irreps retain those indices,
and fermionic irreps follow them. Separate $\alpha/\beta$ counts are not
defined for GHF. ROHF and CUHF require the minority-spin count to be no larger
than the majority-spin count in every irrep.

\section{Molecular frame and point-group detection}
\label{sec:geometry}

\code{GeometryHelper} forms the center of mass and inertia tensor from atomic
masses $m_a$ and positions $\bvec{r}_a$:
\begin{equation}
 \bvec{r}_{\mathrm{COM}}=\frac{\sum_a m_a\bvec{r}_a}{\sum_a m_a},\qquad
 \bvec{I}=\sum_a m_a\left[
  |\bvec{q}_a|^2\Id_3-\bvec{q}_a\bvec{q}_a^{\mathsf T}\right],
 \quad\bvec{q}_a=\bvec{r}_a-\bvec{r}_{\mathrm{COM}}.
 \label{eq:inertia}
\end{equation}
\code{PGSymmetryDetector} diagonalizes $\bvec{I}$ and classifies the molecule
according to degeneracies of its principal moments. For an asymmetric top,
the inertia eigenvectors determine the axes, with axis ordering chosen using
the available twofold rotations. For a symmetric top, the unique inertia
axis determines $z$; candidate perpendicular twofold axes through atoms or
equivalent-atom midpoints, or a vertical mirror plane, fix the transverse
frame. A linear molecule uses its zero-moment axis as $z$, with an arbitrary
orthonormal transverse frame. For spherical tops, candidate twofold axes
and, for octahedral geometry, fourfold axes determine the frame. Cases for
which the supported frame cannot be established, including icosahedral
geometry, use $C_1$.

The rotation matrix is made orthogonal if necessary and adjusted to have
determinant $+1$. Centered coordinates are transformed into this right-handed
frame, and these principal-frame coordinates replace the input coordinates
used by subsequent integral and symmetry operations.

The detector tests inversion, the three Cartesian twofold rotations, and the
three Cartesian mirror planes. An operation is present if every transformed
position coincides, within the geometric tolerance, with an atom of the same
element. The tuple consisting of inversion presence, the number of twofold
axes, and the number of mirror planes selects the group in
\cref{tab:group-detection}. A single atom at the origin is assigned $D_{2h}$
as its supported spatial subgroup.

\begin{table}[ht]
\setlength{\tabcolsep}{10pt}
\caption{Supported point groups selected from the operation tests.}
\label{tab:group-detection}
\begin{tabular}{@{}lccc@{}}
\hline
Group & Inversion & Twofold axes & Mirror planes \\
\hline
$C_1$ & No & 0 & 0 \\
$C_i$ & Yes & 0 & 0 \\
$C_2$ & No & 1 & 0 \\
$C_s$ & No & 0 & 1 \\
$D_2$ & No & 3 & 0 \\
$C_{2v}$ & No & 1 & 2 \\
$C_{2h}$ & Yes & 1 & 1 \\
$D_{2h}$ & Yes & 3 & 3 \\
\hline
\end{tabular}
\end{table}

\section{Spatial symmetry in the AO basis}
\label{sec:ao}

\subsection{Permutation and phase representation}

Define $\bvec{U}(g)$ by its action on AO coefficient vectors:
\begin{equation}
 \hat U(g)\chi_\nu=\sum_\mu\chi_\mu U_{\mu\nu}(g).
 \label{eq:ao-action}
\end{equation}
For the supported spatial groups, each operation maps an atom onto an
equivalent atom and each real spherical harmonic onto itself up to a sign.
Consequently $\bvec{U}(g)$ is a permutation matrix multiplied by local phases.
The builder transforms each atomic position, locates an atom of the same
element at that position, and matches basis functions with the same shell and
angular labels. The radial part of the basis function is unchanged.

\Cref{tab:phases} gives the phases used by \code{local_sign}. Here $l$ is the
angular momentum and $q=|m|$; positive and negative $m$ label the cosine and
sine real harmonics, respectively. The $m=0$ column is included in $m\geq0$.
\begin{table}[ht]
\setlength{\tabcolsep}{10pt}
\caption{Local real-spherical-harmonic phases in the principal frame.}
\label{tab:phases}
\begin{tabular}{@{}lll@{}}
\hline
Operation & $m\geq0$ & $m<0$ \\
\hline
$E$ & $1$ & $1$ \\
$i$ & $(-1)^l$ & $(-1)^l$ \\
$C_{2z}$ & $(-1)^q$ & $(-1)^q$ \\
$C_{2x}$ & $(-1)^{l+q}$ & $(-1)^{l+q+1}$ \\
$C_{2y}$ & $(-1)^l$ & $(-1)^{l+1}$ \\
$\sigma_{xz}$ & $1$ & $-1$ \\
$\sigma_{yz}$ & $(-1)^q$ & $(-1)^{q+1}$ \\
$\sigma_{xy}$ & $(-1)^{l+q}$ & $(-1)^{l+q}$ \\
\hline
\end{tabular}
\end{table}

The coefficient transformation is distinct from the AO matrix element of the
operator. Using \cref{eq:ao-action},
\begin{equation}
 \braket{\chi_\mu|\hat U(g)|\chi_\nu}
   =\bigl[\bvec{S}\bvec{U}(g)\bigr]_{\mu\nu},\qquad
 \bvec{R}(g)=\bvec{C}^\dagger\bvec{S}\bvec{U}(g)\bvec{C}.
 \label{eq:mo-rep}
\end{equation}
The overlap metric is essential in the MO representation. In a closed AO
space,
\begin{equation}
 \bvec{U}(g)^\dagger\bvec{S}\bvec{U}(g)=\bvec{S}.
 \label{eq:metric-invariance}
\end{equation}
Consequently $\bvec{R}(g)$ is unitary whenever the MO space is invariant under
$g$.

\subsection{Spatial character labels and products}

The spatial character tables contain only one-dimensional real irreps with
characters $\pm1$. A symmetry-adapted orbital therefore satisfies
\begin{equation}
 \hat U(g)\varphi_p=\chi_{\Gamma_p}(g)\varphi_p,
 \qquad R_{pq}(g)=\delta_{pq}\chi_{\Gamma_p}(g).
 \label{eq:adapted}
\end{equation}
The orbital's character vector across all operations is matched to the table.
The Cotton index convention makes direct products an exclusive-or operation,
\begin{equation}
 h(\Gamma\otimes\Lambda)=h(\Gamma)\mathbin{\mathrm{XOR}}h(\Lambda).
 \label{eq:xor}
\end{equation}
For example, the $D_{2h}$ ordering is
$a_g,b_{1g},b_{2g},b_{3g},a_u,b_{1u},b_{2u},b_{3u}$.
The determinant irrep is the product of its occupied spin-orbital irreps.
Two occurrences of any spatial irrep cancel in \cref{eq:xor}, so doubly
occupied spatial orbitals contribute the identity. Every closed-shell RHF
determinant is therefore totally symmetric.

\section{Resolving and preserving orbital symmetry}
\label{sec:spatial-basis}

\subsection{Invariant MO blocks}

\code{MOSymmetryDetector} computes all representations in
\cref{eq:mo-rep}. A graph connects MO indices $p$ and $q$ if any operation has
$|R_{pq}(g)|$ above the representation tolerance. The adjacency is symmetrized,
and each connected component defines a candidate invariant block. The
construction includes nonadjacent MO columns and is independent of orbital
energy differences. Energy proximity alone does not establish closure under
the symmetry operations.

Within each connected block, the detector simultaneously resolves the
commuting operations. Starting with the whole block, it diagonalizes one
operation and separates its positive and negative eigenspaces. The next
operation is diagonalized only within each of these subspaces. Repeating this
process preserves every previously resolved character while subdividing the
space into irreps.

Let $\bvec{V}_\Gamma$ contain the resulting rotation vectors in one irrep
subspace and let $\bvec{e}=\operatorname{diag}(\epsilon_p)$ be the supplied
orbital-energy operator in the original block. The detector diagonalizes
\begin{equation}
 \bvec{e}_\Gamma=\bvec{V}_\Gamma^\dagger\bvec{e}\bvec{V}_\Gamma
 \label{eq:energy-proxy}
\end{equation}
to fix the remaining freedom among repeated copies of the same irrep. The
updated coefficients and projected energies are sorted together within the
connected block. If the input orbitals are not eigenvectors of a
symmetry-invariant operator, this is a projected energy convention rather
than preservation of each input orbital energy.

After adaptation, every operation must have a diagonal representation of unit
magnitude, and each character vector must match a table row. The default
absolute tolerance is $10^{-6}$, with zero relative tolerance. Failure raises
\code{RuntimeError}; an unresolved orbital is not assigned an arbitrary label.
The detector updates both coefficients and energies in place.

\subsection{Symmetry-block Fock diagonalization}

\code{SymmetryBasis.build} applies the spatial detector to $\bvec{X}$ with
zero energy proxies, producing orthonormal AO coefficient blocks
$\bvec{B}_\Gamma$. For any Hermitian AO Fock matrix $\bvec{F}$, it solves
\begin{equation}
 \bvec{F}_\Gamma=\bvec{B}_\Gamma^\dagger\bvec{F}\bvec{B}_\Gamma,
 \qquad
 \bvec{F}_\Gamma\bvec{v}_\Gamma
   =\bvec{v}_\Gamma\bvec{\epsilon}_\Gamma,
 \qquad
 \bvec{C}_\Gamma=\bvec{B}_\Gamma\bvec{v}_\Gamma.
 \label{eq:block-fock}
\end{equation}
The complete result is stably sorted by energy, carrying coefficients and
irrep indices through the same permutation. For an invariant Fock operator
the off-diagonal irrep blocks vanish; restricting diagonalization to the
diagonal blocks also prevents numerical or guess-induced inter-irrep mixing.

Spin-free GHF uses the spatial action $\Id_2\otimes\bvec{U}(g)$ in the
two-component AO basis. This duplicates the spatial representation without
rotating the spin. Spin--orbit GHF instead uses the double-group construction
in \cref{sec:double-group}.

\section{Occupation selection}
\label{sec:selection}

An occupation policy returns a permutation of MO columns. The same
permutation is applied to coefficients, energies, and irrep metadata before
the density is constructed. Occupied orbitals precede virtual orbitals, with
stable energy ordering within each partition. ROHF additionally separates
doubly occupied, singly occupied, and virtual columns. These conventions
preserve the occupied-column slices used by the density builders.

\subsection{Explicit irrep counts}

With explicit counts, the policy takes the lowest-energy orbitals within each
irrep. It checks that counts are nonnegative integers, their totals match the
electron numbers, and no count exceeds the corresponding basis capacity.
Equivalent label and integer keys cannot specify an irrep twice. If a target
is also supplied, the configuration must realize it. ROHF and CUHF enforce
nested minority/majority counts in every irrep.

\subsection{A target in a one-dimensional irrep algebra}

Target-only selection minimizes the sum of occupied orbital energies at the
current iteration. This local occupation criterion does not guarantee the
globally lowest HF stationary solution in the target irrep.

For $M$ orbitals and a requested occupation $N$, let $A_p(k,h)$ be the minimum
energy sum obtained by occupying $k$ of the first $p$ orbitals with total irrep
$h$. For orbital $p$ with energy $\epsilon_p$ and irrep $\gamma_p$,
\begin{equation}
 A_p(k,h)=\min\left\{
 A_{p-1}(k,h),\,
 A_{p-1}(k-1,h\otimes\gamma_p^{-1})+\epsilon_p
 \right\}.
 \label{eq:dp}
\end{equation}
The initial condition is $A_0(0,0)=0$, with all other entries infinite.
Backtracking from $A_M(N,h_{\mathrm{target}})$ yields the occupied indices.
The algorithm and its stored backtracking history cost $O(MNR)$ for $R$
irreps. Spatial products use \cref{eq:xor}; Abelian double-group products use
the character-derived multiplication table in \cref{eq:products}.

UHF constructs independent energy spectra $A_\alpha(h)$ and $A_\beta(h)$ at
the required spin electron counts, then minimizes
\begin{equation}
 \min_h\left[
 A_\alpha(h)+A_\beta(h^{-1}\otimes h_{\mathrm{target}})
 \right].
 \label{eq:uhf-selection}
\end{equation}
Thus both spins are selected jointly subject to their combined target.

For ROHF and CUHF the policy uses a dynamic program over core count, open-shell
count, and total irrep. In each irrep, let $c$ be the minority-spin count and
$s$ the number of additional majority-spin orbitals. With $P_\sigma(n)$ the
prefix sum of the lowest $n$ orbital energies of spin $\sigma$, the local cost
is
\begin{equation}
 E_\Gamma(c,s)=P_{\mathrm{minor}}(c)
              +P_{\mathrm{major}}(c+s).
 \label{eq:nested-cost}
\end{equation}
ROHF uses its common orbital spectrum for both prefix sums. The transition
adds $c$ core and $s$ open-shell occupations and contributes
$\Gamma^{\otimes s}$ to the total irrep. The final counts must be
$N_{\mathrm{core}}=\min(N_\alpha,N_\beta)$ and
$N_{\mathrm{open}}=|N_\alpha-N_\beta|$. RHF permits only an identity target;
explicit pair counts can select among closed-shell configurations sharing it.

\section{Double groups for spin--orbit GHF}
\label{sec:double-group}

\subsection{Spin lifts and character tables}

Two-component coefficients use the block layout
\code{[alpha AOs, beta AOs]}. The unbarred lift of a spatial operation is
\begin{equation}
 \bvec{U}^*(g)=\bvec{D}^{1/2}(g)\otimes\bvec{U}(g),\qquad
 \bvec{D}^{1/2}(C_{2k})=-\mathrm{i}\bvec{\sigma}_k,
 \label{eq:spin-lift}
\end{equation}
where $\bvec{\sigma}_k$ is a Pauli matrix. Both identity and inversion act as
$\Id_2$ on spin. A mirror is inversion times the $\pi$ rotation about its
normal and therefore uses the same spin matrix as that rotation.

The double group $G^*$ contains an unbarred and a barred lift for every spatial
operation. The central $2\pi$ rotation, $\bar E$, acts as $-\Id$ on a spinor,
and
\begin{equation}
 \bvec{U}^*(\bar g)=-\bvec{U}^*(g).
 \label{eq:barred}
\end{equation}
Bosonic characters are unchanged by barring, while fermionic characters
change sign. \code{DoubleGroup} stores all unbarred operations followed by
all barred operations in the same spatial order.

\Cref{tab:spinors} lists the fermionic labels. In $C_2^*$, $C_s^*$, and
$C_{2h}^*$, the labels beginning with \code{1} and \code{2} have characters
$+\mathrm{i}$ and $-\mathrm{i}$ under the unbarred $C_2$ rotation (the mirror
for $C_s^*$). Their products are not given by spatial Cotton-index XOR.
\begin{table}[ht]
\setlength{\tabcolsep}{10pt}
\caption{Fermionic irreps appended after the ordinary Cotton-indexed irreps.}
\label{tab:spinors}
\begin{tabular}{@{}lll@{}}
\hline
Double group & Implementation labels & Dimension \\
\hline
$C_1^*$ & \code{a1/2} & 1 \\
$C_i^*$ & \code{a1/2g}, \code{a1/2u} & 1 \\
$C_2^*$, $C_s^*$ & \code{1e1/2}, \code{2e1/2} & 1 \\
$C_{2h}^*$ & \code{1e1/2g}, \code{2e1/2g}, & 1 \\
 & \code{1e1/2u}, \code{2e1/2u} & \\
$D_2^*$, $C_{2v}^*$ & \code{e1/2} & 2 \\
$D_{2h}^*$ & \code{e1/2g}, \code{e1/2u} & 2 \\
\hline
\end{tabular}
\end{table}

Direct-product multiplicities are computed from the characters:
\begin{equation}
 N_{\Gamma\Lambda}^{\Omega}
 =\frac{1}{|G^*|}\sum_{g\in G^*}
   \chi_\Gamma(g)\chi_\Lambda(g)\chi_\Omega(g)^*.
 \label{eq:products}
\end{equation}
For an Abelian double group exactly one multiplicity is one, giving the
product table used in \cref{eq:dp}. For the non-Abelian groups products can be
reducible. For example,
\begin{equation}
 E_{1/2,g}\otimes E_{1/2,g}
   =A_g\oplus B_{1g}\oplus B_{2g}\oplus B_{3g}
 \quad\text{in }D_{2h}^*.
 \label{eq:reducible}
\end{equation}
An occupied-spinor label product therefore does not by itself determine the
total irrep of a Slater determinant.

\subsection{Orthonormal spinor blocks}

Noncommuting spin rotations cannot be simultaneously diagonalized as in the
spatial detector. \code{DoubleGroupBasis} instead constructs character
projectors in the retained orthonormal AO space:
\begin{equation}
 \bvec{P}_\Gamma
  =\frac{d_\Gamma}{|G^*|}\sum_{g\in G^*}\chi_\Gamma(g)^*
    \bvec{X}^\dagger\bvec{S}\bvec{U}^*(g)\bvec{X}.
 \label{eq:spinor-projector}
\end{equation}
The eigenvectors with eigenvalue one span the isotypic block for $\Gamma$,
including every equivalent copy of that irrep. Multiplying these vectors by
$\bvec{X}$ produces the AO coefficient block $\bvec{B}_\Gamma$ used in
\cref{eq:block-fock}.

Before building these blocks, the implementation verifies that every
representation $\bvec{X}^\dagger\bvec{S}\bvec{U}^*(g)\bvec{X}$ is unitary.
It also checks that projector eigenvalues are zero or one and that the blocks
span the full retained space with the correct overlap orthonormality. These
checks reject, for example, an orthogonalizer that retains only one component
of a spinor multiplet. The absolute tolerance is $10^{-6}$ with zero relative
tolerance.

A spinor labelled by a two-dimensional irrep belongs to its isotypic block;
an individual column need not be an eigenvector of every operation.
Unpaired block diagonalization preserves this membership while allowing the
Fock eigenvectors to mix equivalent copies.

\section{Total determinant symmetry and allowed targets}
\label{sec:determinant}

\subsection{Slater determinant projectors}

Let $\bvec{C}_{\mathrm{occ}}$ contain the $N$ occupied orthonormal spinors.
The overlap with the symmetry-transformed determinant is
\begin{equation}
 f(g)=\braket{\Phi|\hat U^*(g)|\Phi}
     =\det\left(
       \bvec{C}_{\mathrm{occ}}^\dagger\bvec{S}
       \bvec{U}^*(g)\bvec{C}_{\mathrm{occ}}
       \right).
 \label{eq:det-overlap}
\end{equation}
The total-irrep weight is the expectation value of its character projector:
\begin{equation}
 w_\Gamma=\frac{d_\Gamma}{|G^*|}
   \sum_{g\in G^*}\chi_\Gamma(g)^*f(g).
 \label{eq:det-weight}
\end{equation}
The weights are real, nonnegative, and sum to one for a normalized
determinant. The implementation verifies nonnegativity and normalization
within $10^{-6}$, clips small numerical excursions to $[0,1]$, and reports a
pure irrep when exactly one weight is one within that tolerance.

\code{hf.state_symmetry_weights} stores the weights by irrep label.
\code{hf.state_symmetry} is the pure label or \code{None} for a mixture. For a
multidimensional irrep, a pure weight means membership in that total isotypic
space; it does not mean the state transforms by a scalar under every
operation.

The central rotation satisfies
$\hat U^*(\bar E)\ket{\Phi}=(-1)^N\ket{\Phi}$.
Thus odd-electron targets must be fermionic and even-electron targets bosonic.
Every requested target is checked against the final determinant through
\cref{eq:det-weight}.

\subsection{Non-Abelian single-determinant restrictions}

For $D_2^*$, $C_{2v}^*$, and $D_{2h}^*$, the one-electron fermionic irreps are
two dimensional. An even-electron determinant transforming as a
one-dimensional total irrep must have an occupied subspace invariant under
every group operation: a transformed determinant proportional to the
original has the same occupied space. Such a space is a direct sum of complete
copies of the spinor irreps.

Each complete copy has determinant one for every operation. Proper spin
rotations have determinant one, the barred lift has determinant
$\det(-\Id_2)=1$, and inversion contributes $\det(\pm\Id_2)=1$. Hence the
determinant of any invariant occupied space is totally symmetric. Other
bosonic targets in these three double groups require a multideterminant
wavefunction and are rejected by the HF occupation policy.

For an odd-electron target in these groups, the fermionic isotypic space is
fixed by inversion parity when inversion is present. The occupation search
then reduces to the parity version of \cref{eq:dp}. $D_2^*$ and $C_{2v}^*$
have only one fermionic irrep, while $D_{2h}^*$ distinguishes its $g$ and $u$
irreps. This reduction follows from the character table of these particular
groups; it is not a general replacement for double-group symmetry.

Explicit counts without a target may describe an even-electron determinant
with several nonzero total-irrep weights. Those counts constrain spinor block
membership, while \cref{eq:det-weight} determines the resulting electronic
mixture. Abelian double groups use the full one-dimensional product table and
do not have the complete-two-dimensional-multiplet restriction.

\subsection{Paired diagonalization for the identity target}
\label{sec:paired}

To realize a totally symmetric even-electron determinant in a non-Abelian
double group, the basis must supply complete partner pairs. Within an
isotypic block, \code{DoubleGroupBasis} constructs and caches two component
bases, $\bvec{B}_0$ and $\bvec{B}_1$. The first spans the $-\mathrm{i}$
eigenspace of $C_{2z}$, and the second is obtained by applying $C_{2y}$ for
$D_2^*$ and $D_{2h}^*$, or $\sigma_{xz}$ for $C_{2v}^*$.

Schur's lemma gives the invariant Fock operator by tracing over the irrep
components. Its multiplicity-space matrix is
\begin{equation}
 \bvec{F}_{\mathrm{eff}}=\frac12\left(
     \bvec{B}_0^\dagger\bvec{F}\bvec{B}_0
    +\bvec{B}_1^\dagger\bvec{F}\bvec{B}_1\right).
 \label{eq:schur}
\end{equation}
If $\bvec{F}_{\mathrm{eff}}\bvec{v}=\bvec{v}\bvec{\epsilon}$, the two
partners are $\bvec{B}_0\bvec{v}$ and $\bvec{B}_1\bvec{v}$ with identical
energies. This component trace is equivalent to the full group average in
the isotypic block,
\begin{equation}
 \overline{\bvec{F}}_\Gamma
 =\frac{1}{|G^*|}\sum_{g\in G^*}
    \bvec{R}_\Gamma(g)^\dagger\bvec{F}_\Gamma\bvec{R}_\Gamma(g).
 \label{eq:group-average}
\end{equation}
Caching the component bases avoids constructing all transformed Fock blocks
in each SCF iteration.

Partners are interleaved before stable global energy sorting. The occupation
policy selects whole pairs, including at accidental energy degeneracies.
With explicit counts and this target, every spinor-irrep count must be even.
Without explicit counts, the lowest-energy pairs are filled.

\section{Shared SCF interfaces and data flow}
\label{sec:implementation}

The spatial and double-group basis classes share two operations:
\begin{itemize}
 \item \code{eigh(F)} returns energies, coefficients, and irrep indices
       from a symmetry-block diagonalization.
 \item \code{adapt(C, eps)} returns the same three quantities for a supplied
       orbital guess.
\end{itemize}
Spatial adaptation uses the cached AO operations and the MO detector. Spinor
adaptation constructs the AO energy proxy
\begin{equation}
 \bvec{F}_{\mathrm{guess}}
 =\bvec{S}\bvec{C}\operatorname{diag}(\epsilon_p)
    \bvec{C}^\dagger\bvec{S},
 \label{eq:guess-proxy}
\end{equation}
then calls the double-group basis diagonalizer. This also enforces paired
structure when required by \cref{sec:paired}. If no guess energies are stored,
the diagonal of $\bvec{C}^\dagger\bvec{H}\bvec{C}$ supplies the proxy, where
$\bvec{H}$ is the core Hamiltonian.

\code{SCFBase} owns the shared orchestration:
\begin{enumerate}
 \item Constructor validation checks option types and count shapes.
 \item Binding resolves the group, electron numbers, labels, and compatible
       occupation constraints.
 \item Run startup constructs the symmetry basis once from $\bvec{X}$ and
       checks capacity and target feasibility before the first density or
       J/K build. Feasibility uses the occupation policy with zero costs.
 \item With constraints active, SAP/core guesses use symmetry-block
       diagonalization, and supplied or perturbed guesses are adapted before
       selecting their occupations.
 \item Each SCF Fock diagonalization returns irrep metadata with its
       coefficients; the occupation policy permutes all three arrays before
       density construction.
 \item Final labels use that metadata. The total determinant irrep is
       computed separately and verified against the requested target.
\end{enumerate}
Spatial symmetry blocks and spin--orbit double-group blocks are also used
when no occupation constraint is requested. In that case ordinary energy
ordering determines occupations and the final determinant symmetry is
reported. Spatial $C_1$ can use the ordinary orthogonalized diagonalizer.

Raw symmetry and occupation options are preserved. Resolved count arrays,
normalized labels, symmetry bases, and cached multiplet components are
derived state. Rebuilding a method chain can therefore pass the original
options unchanged to a new constructor. Rebinding clears orbital and
determinant metadata, and the basis caches are reconstructed at run startup.

The source responsibilities are:
\begin{description}
 \item[\code{system/geom_utils.py}] Center of mass, inertia tensor, and
       principal-frame coordinates.
 \item[\code{symmetry/pg_sym_detect.py}] Frame selection and molecular
       operation tests for the supported spatial point groups.
 \item[\code{symmetry/sym_utils.py}] Operation order, spatial character
       tables, Cotton indices, and geometric operations.
 \item[\code{symmetry/mo_sym_detect.py}] AO permutation/phase matrices,
       connected MO blocks, simultaneous spatial resolution, and character
       validation.
 \item[\code{symmetry/symmetry_basis.py}] Spatial orthonormal irrep blocks
       and the common diagonalization/adaptation interface.
 \item[\code{symmetry/double_group.py}] Spin lifts, character tables and
       products, spinor projectors, paired blocks, and determinant weights.
 \item[\code{scf/occupations.py}] Shared input/capacity validation and
       separate spatial and double-group occupation policies.
 \item[\code{scf/scf_base.py}] SCF lifecycle, guess preparation, consistent
       orbital permutations, and final determinant verification.
\end{description}
The two occupation policies share \code{_OccupationConstraints} for label
resolution and capacity checks. The double-group policy is independent of the
spatial policy, so shared validation carries no spatial XOR assumption.

\section{Boundary with relativistic CI}
\label{sec:ci}

The current CI string algebra implements one-dimensional spatial XOR
products. Spin--orbit \code{RelCIBase} therefore supplies $C_1$ orbital indices
to its CI workers and requires the CI state-symmetry index to be zero.
\code{CIBase._active_orbital_symmetries} provides the hook for this conversion.
Spin-free relativistic CI retains spatial symmetry handling.

An HF double-group target constrains the HF reference determinant. It does
not select a subsequent CI or MCSCF electronic irrep. Full double-group
state selection in those methods requires a determinant/state representation
that handles the reducible products and multidimensional irreps explicitly.

\end{document}
